\documentclass[hidelinks]{article}
\usepackage{stex}

\libinput{preambleCS}

\title{Type 0 and Type 1 Languages and their Automata}
\author{}
\date{}

\begin{document}

\maketitle

\begin{smodule}[title={Type 0 and Type 1 Languages and their Automata}]{Type0Type1Languages}
\usemodule{ComputerScience/FormalLanguagesAndAutomata?IntroductionToFormalLanguagesAndAutomata}
\usemodule{ComputerScience/FormalLanguagesAndAutomata?StandardTuringMachines}

\symdecl*{recursively enumerable language}
\symdecl*{recursive language}
\symdecl*{unrestricted grammar}
\symdecl*{context-sensitive grammar}
\symdecl*{context-sensitive language}
\symdecl*{linear-bounded automaton}

\section{Overview}

This note populates the top two levels of the Chomsky hierarchy. The language class accepted by \sns{Turing machine} is the family of \sns{recursively enumerable language} (Type 0); the corresponding grammar class is the family of \sns{unrestricted grammar}. Restricting the grammars so that no production reduces the working string in length yields the \sns{context-sensitive grammar} (Type 1) and \sns{context-sensitive language}, which are recognised by the \sns{linear-bounded automaton}.

\section{Recursive and recursively enumerable languages}

\begin{sdefinition}[for=recursively enumerable language]
A \sn{formal language} $L$ is a \definame{recursively enumerable language} (Type 0) iff there is a \sn{Turing machine} that accepts it. Strings $w \notin L$ may be rejected because the machine halts in a non-final state \emph{or} because the machine fails to halt at all.
\end{sdefinition}

\begin{sdefinition}[for=recursive language]
A \sn{formal language} $L$ is a \definame{recursive language} iff there is a \sn{Turing machine} that accepts it which is guaranteed to halt on every input string. Equivalently, $L$ is recursive iff it has a \emph{membership algorithm}.
\end{sdefinition}

\subsection{Non-recursive and non-recursively-enumerable languages}

\begin{sparagraph}[title={Theorem: not every language is recursively enumerable}]
For any non-empty alphabet $\Sigma$, there are languages over $\Sigma$ that are not \sn{recursively enumerable language}.
\end{sparagraph}

\begin{sparagraph}[title={Proof sketch}]
A diagonalisation argument, based on the fact that there are uncountably many subsets of $\Sigma^*$ but only countably many \sns{Turing machine}, gives the existence of a non-RE language. Linz and other textbooks contain the complete argument. The takeaway is that there are more languages than Turing machines.
\end{sparagraph}

\begin{sparagraph}[title={Theorem: recursive $\subsetneq$ recursively enumerable}]
The \sns{recursive language} are a strict subset of the \sns{recursively enumerable language}.
\end{sparagraph}

\section{Unrestricted grammars}

\begin{sdefinition}[for=unrestricted grammar]
An \definame{unrestricted grammar} (sometimes called a Type-0 grammar) is a phrase-structure grammar with no restrictions on its productions, except that the empty string $\lambda$ may not appear on the left-hand side. Productions have the general form $\alpha \to \beta$ in which $\alpha$ and $\beta$ are arbitrary strings of terminals and non-terminals.
\end{sdefinition}

\begin{example}[An unrestricted grammar]
The following productions are unrestricted:
\[
\begin{array}{rcl}
S &\to& S_1 B,\\
S_1 &\to& a S_1 b,\\
b B &\to& b b b B,\\
a S_1 b &\to& a a,\\
B &\to& \lambda.
\end{array}
\]
A short derivation analysis shows that the language generated is $\{a^{n+1} b^{n-1} b^{2m} \mid n \geq 1,\ m \geq 0\}$.
\end{example}

\subsection{Unrestricted grammars and recursively enumerable languages}

Every language generated by an \sn{unrestricted grammar} is a \sn{recursively enumerable language}: enumerate strings of terminals by length of derivation. Conversely, every \sn{recursively enumerable language} (accepted by some TM) can be generated by an \sn{unrestricted grammar}; the construction is complex and lengthy, but uses essentially the same idea as the NPDA-to-CFG construction. So the \sns{unrestricted grammar} generate exactly the \sns{recursively enumerable language}, which is the largest family of languages that can be generated or recognised algorithmically.

\section{Context-sensitive grammars}

\begin{sdefinition}[for=context-sensitive grammar]
A \definame{context-sensitive grammar} (Type 1) is a phrase-structure grammar in which the right-hand side of every production is at least as long as the left-hand side. Equivalently, every production has the form $\alpha A \gamma \to \alpha \beta \gamma$ with $\beta \neq \lambda$, so that the non-terminal $A$ is replaced by $\beta$ only in the context of $\alpha$ and $\gamma$.
\end{sdefinition}

\begin{example}[A context-sensitive grammar for $a^n b^n c^n$]
The grammar
\[
\begin{array}{rcl}
S &\to& abc \mid aAbc,\\
Ab &\to& bA,\\
Ac &\to& Bbcc,\\
bB &\to& Bb,\\
aB &\to& aa \mid aaA
\end{array}
\]
is context-sensitive (each production has $|\beta| \geq |\alpha|$). It generates the \sn{context-sensitive language} $L = \{a^n b^n c^n \mid n \geq 1\}$. A typical derivation is
\[
S \Rightarrow a\underline{A}bc \Rightarrow ab\underline{Ac} \Rightarrow a\underline{bB}bcc \Rightarrow a\underline{aB}bcc \Rightarrow aabbcc \Rightarrow \cdots
\]
With the non-terminals $A$ and $B$ acting as messengers, moving through the string as it develops, this gives the canonical example showing that context-sensitive languages are strictly more expressive than context-free languages.
\end{example}

\subsection{Properties}

\sns{context-sensitive grammar} have two characteristic properties:
\begin{itemize}
  \item \emph{non-contracting}: the working string of terminals and non-terminals cannot reduce in length at any point;
  \item \emph{context-sensitive}: a non-terminal can only be replaced in certain contexts (in the example above, $A$ should be replaced only when followed by $b$ or $c$).
\end{itemize}

\subsection{Context-sensitive languages}

\begin{sdefinition}[for=context-sensitive language]
A \sn{formal language} $L$ is a \definame{context-sensitive language} (Type 1) iff there is a \sn{context-sensitive grammar} $G$ such that $L = L(G)$ or $L = L(G) \cup \{\lambda\}$.
\end{sdefinition}

The empty string $\lambda$ is included by definition because a \sn{context-sensitive grammar} can never generate $\lambda$ (the productions are non-contracting). The family of context-free languages is a strict subset of the family of \sns{context-sensitive language}.

\section{Linear bounded automata}

\begin{sdefinition}[for=linear-bounded automaton]
A \definame{linear-bounded automaton} (LBA) is similar to a standard \sn{Turing machine}, but with the restriction that it can only use the part of the tape that originally held the input. One way to think about this is that the ends of the input are marked with special markers, say $[$ and $]$, that cannot be moved, passed over, or overwritten.
\end{sdefinition}

For example, the standard TM described previously for $L = \{a^n b^n\}$ examines and overwrites the input symbols but never uses extra tape space, so it is in fact a \sn{linear-bounded automaton}. A similar idea works to recognise $\{a^n b^n c^n\}$.

\subsection{Equivalence with context-sensitive grammars}

\begin{sparagraph}[title={Theorem: every CSL is accepted by an LBA}]
For every \sn{context-sensitive language} not including $\lambda$, there exists a \sn{linear-bounded automaton} that accepts it.
\end{sparagraph}

\begin{sparagraph}[title={Theorem: every LBA-accepted language is a CSL}]
If $L$ is accepted by a \sn{linear-bounded automaton}, then there is a \sn{context-sensitive grammar} that generates $L$.
\end{sparagraph}

So \sns{context-sensitive grammar} generate exactly the \sns{context-sensitive language}, i.e.\ those languages accepted by \sns{linear-bounded automaton}.

\subsection{Recursive vs context-sensitive}

Every \sn{context-sensitive language} is a \sn{recursive language} (it has a membership algorithm: enumerate derivations up to length bounded by the input). However, some \sns{recursive language} are not context-sensitive, so the inclusion is strict. Ordering the expressive power of automata models, a \sn{linear-bounded automaton} is \emph{less} powerful than a \sn{Turing machine} but \emph{more} powerful than a non-deterministic pushdown automaton.

\section{The Chomsky hierarchy revisited}

The four families introduced across this course form the Chomsky hierarchy:
\[
\begin{array}{|c|l|l|l|}
\hline
\text{Type} & \text{Grammar/Language} & \text{Grammar Constraint} & \text{Machine}\\
\hline
0 & \text{Recursively enumerable} & \alpha \to \beta & \text{Turing machines}\\
1 & \text{Context-sensitive} & \alpha A \gamma \to \alpha \beta \gamma & \text{Linear-bounded automata}\\
2 & \text{Context-free} & A \to \beta & \text{Non-deterministic pushdown automata}\\
3 & \text{Regular} & A \to aB & \text{Finite-state automata}\\
\hline
\end{array}
\]
The constraint on Type 1 grammars is the length restriction $|\alpha A \gamma| \leq |\alpha \beta \gamma|$, equivalently $\beta \neq \lambda$.

\section{Summary}

The \sn{recursively enumerable language} (Type 0) are exactly the languages accepted by \sns{Turing machine} and generated by \sns{unrestricted grammar}; some are not recursive (have no membership algorithm) and some are not even recursively enumerable. The \sn{context-sensitive language} (Type 1) are those generated by the non-contracting \sns{context-sensitive grammar} and accepted by \sns{linear-bounded automaton}; they fit strictly between the context-free and recursive language classes. Together with the regular and context-free families from earlier in the course, this completes the Chomsky hierarchy.

\end{smodule}

\end{document}
